\IfFileExists{stacks-project.cls}{%
\documentclass{stacks-project}
}{%
\documentclass{amsart}
}

% For dealing with references we use the comment environment
\usepackage{verbatim}
\newenvironment{reference}{\comment}{\endcomment}
%\newenvironment{reference}{}{}
\newenvironment{slogan}{\comment}{\endcomment}
\newenvironment{history}{\comment}{\endcomment}

% For commutative diagrams we use Xy-pic
\usepackage[all]{xy}

% We use 2cell for 2-commutative diagrams.
\xyoption{2cell}
\UseAllTwocells

% We use multicol for the list of chapters between chapters
\usepackage{multicol}

% This is generally recommended for better output
\usepackage{lmodern}
\usepackage[T1]{fontenc}

% For cross-file-references
\usepackage{xr-hyper}

% Package for hypertext links:
\usepackage{hyperref}

% For any local file, say "hello.tex" you want to link to please
% use \externaldocument[hello-]{hello}
\externaldocument[introduction-]{introduction}
\externaldocument[conventions-]{conventions}
\externaldocument[sets-]{sets}
\externaldocument[categories-]{categories}
\externaldocument[topology-]{topology}
\externaldocument[sheaves-]{sheaves}
\externaldocument[sites-]{sites}
\externaldocument[stacks-]{stacks}
\externaldocument[fields-]{fields}
\externaldocument[algebra-]{algebra}
\externaldocument[brauer-]{brauer}
\externaldocument[homology-]{homology}
\externaldocument[derived-]{derived}
\externaldocument[simplicial-]{simplicial}
\externaldocument[more-algebra-]{more-algebra}
\externaldocument[smoothing-]{smoothing}
\externaldocument[modules-]{modules}
\externaldocument[sites-modules-]{sites-modules}
\externaldocument[injectives-]{injectives}
\externaldocument[cohomology-]{cohomology}
\externaldocument[sites-cohomology-]{sites-cohomology}
\externaldocument[dga-]{dga}
\externaldocument[dpa-]{dpa}
\externaldocument[sdga-]{sdga}
\externaldocument[hypercovering-]{hypercovering}
\externaldocument[schemes-]{schemes}
\externaldocument[constructions-]{constructions}
\externaldocument[properties-]{properties}
\externaldocument[morphisms-]{morphisms}
\externaldocument[coherent-]{coherent}
\externaldocument[divisors-]{divisors}
\externaldocument[limits-]{limits}
\externaldocument[varieties-]{varieties}
\externaldocument[topologies-]{topologies}
\externaldocument[descent-]{descent}
\externaldocument[perfect-]{perfect}
\externaldocument[more-morphisms-]{more-morphisms}
\externaldocument[flat-]{flat}
\externaldocument[groupoids-]{groupoids}
\externaldocument[more-groupoids-]{more-groupoids}
\externaldocument[etale-]{etale}
\externaldocument[chow-]{chow}
\externaldocument[intersection-]{intersection}
\externaldocument[pic-]{pic}
\externaldocument[weil-]{weil}
\externaldocument[adequate-]{adequate}
\externaldocument[dualizing-]{dualizing}
\externaldocument[duality-]{duality}
\externaldocument[discriminant-]{discriminant}
\externaldocument[derham-]{derham}
\externaldocument[local-cohomology-]{local-cohomology}
\externaldocument[algebraization-]{algebraization}
\externaldocument[curves-]{curves}
\externaldocument[resolve-]{resolve}
\externaldocument[models-]{models}
\externaldocument[functors-]{functors}
\externaldocument[equiv-]{equiv}
\externaldocument[pione-]{pione}
\externaldocument[etale-cohomology-]{etale-cohomology}
\externaldocument[proetale-]{proetale}
\externaldocument[relative-cycles-]{relative-cycles}
\externaldocument[more-etale-]{more-etale}
\externaldocument[trace-]{trace}
\externaldocument[crystalline-]{crystalline}
\externaldocument[spaces-]{spaces}
\externaldocument[spaces-properties-]{spaces-properties}
\externaldocument[spaces-morphisms-]{spaces-morphisms}
\externaldocument[decent-spaces-]{decent-spaces}
\externaldocument[spaces-cohomology-]{spaces-cohomology}
\externaldocument[spaces-limits-]{spaces-limits}
\externaldocument[spaces-divisors-]{spaces-divisors}
\externaldocument[spaces-over-fields-]{spaces-over-fields}
\externaldocument[spaces-topologies-]{spaces-topologies}
\externaldocument[spaces-descent-]{spaces-descent}
\externaldocument[spaces-perfect-]{spaces-perfect}
\externaldocument[spaces-more-morphisms-]{spaces-more-morphisms}
\externaldocument[spaces-flat-]{spaces-flat}
\externaldocument[spaces-groupoids-]{spaces-groupoids}
\externaldocument[spaces-more-groupoids-]{spaces-more-groupoids}
\externaldocument[bootstrap-]{bootstrap}
\externaldocument[spaces-pushouts-]{spaces-pushouts}
\externaldocument[spaces-chow-]{spaces-chow}
\externaldocument[groupoids-quotients-]{groupoids-quotients}
\externaldocument[spaces-more-cohomology-]{spaces-more-cohomology}
\externaldocument[spaces-simplicial-]{spaces-simplicial}
\externaldocument[spaces-duality-]{spaces-duality}
\externaldocument[formal-spaces-]{formal-spaces}
\externaldocument[restricted-]{restricted}
\externaldocument[spaces-resolve-]{spaces-resolve}
\externaldocument[formal-defos-]{formal-defos}
\externaldocument[defos-]{defos}
\externaldocument[cotangent-]{cotangent}
\externaldocument[examples-defos-]{examples-defos}
\externaldocument[algebraic-]{algebraic}
\externaldocument[examples-stacks-]{examples-stacks}
\externaldocument[stacks-sheaves-]{stacks-sheaves}
\externaldocument[criteria-]{criteria}
\externaldocument[artin-]{artin}
\externaldocument[quot-]{quot}
\externaldocument[stacks-properties-]{stacks-properties}
\externaldocument[stacks-morphisms-]{stacks-morphisms}
\externaldocument[stacks-limits-]{stacks-limits}
\externaldocument[stacks-cohomology-]{stacks-cohomology}
\externaldocument[stacks-perfect-]{stacks-perfect}
\externaldocument[stacks-introduction-]{stacks-introduction}
\externaldocument[stacks-more-morphisms-]{stacks-more-morphisms}
\externaldocument[stacks-geometry-]{stacks-geometry}
\externaldocument[moduli-]{moduli}
\externaldocument[moduli-curves-]{moduli-curves}
\externaldocument[examples-]{examples}
\externaldocument[exercises-]{exercises}
\externaldocument[guide-]{guide}
\externaldocument[desirables-]{desirables}
\externaldocument[coding-]{coding}
\externaldocument[obsolete-]{obsolete}
\externaldocument[fdl-]{fdl}
\externaldocument[index-]{index}

% Theorem environments.
%
\theoremstyle{plain}
\newtheorem{theorem}[subsection]{Theorem}
\newtheorem{proposition}[subsection]{Proposition}
\newtheorem{lemma}[subsection]{Lemma}

\theoremstyle{definition}
\newtheorem{definition}[subsection]{Definition}
\newtheorem{example}[subsection]{Example}
\newtheorem{exercise}[subsection]{Exercise}
\newtheorem{situation}[subsection]{Situation}

\theoremstyle{remark}
\newtheorem{remark}[subsection]{Remark}
\newtheorem{remarks}[subsection]{Remarks}

\numberwithin{equation}{subsection}

% Macros
%
\def\lim{\mathop{\mathrm{lim}}\nolimits}
\def\colim{\mathop{\mathrm{colim}}\nolimits}
\def\Spec{\mathop{\mathrm{Spec}}}
\def\Hom{\mathop{\mathrm{Hom}}\nolimits}
\def\Ext{\mathop{\mathrm{Ext}}\nolimits}
\def\SheafHom{\mathop{\mathcal{H}\!\mathit{om}}\nolimits}
\def\SheafExt{\mathop{\mathcal{E}\!\mathit{xt}}\nolimits}
\def\Sch{\mathit{Sch}}
\def\Mor{\mathop{\mathrm{Mor}}\nolimits}
\def\Ob{\mathop{\mathrm{Ob}}\nolimits}
\def\Sh{\mathop{\mathit{Sh}}\nolimits}
\def\NL{\mathop{N\!L}\nolimits}
\def\CH{\mathop{\mathrm{CH}}\nolimits}
\def\proetale{{pro\text{-}\acute{e}tale}}
\def\etale{{\acute{e}tale}}
\def\QCoh{\mathit{QCoh}}
\def\Ker{\mathop{\mathrm{Ker}}}
\def\Im{\mathop{\mathrm{Im}}}
\def\Coker{\mathop{\mathrm{Coker}}}
\def\Coim{\mathop{\mathrm{Coim}}}

% Boxtimes
%
\DeclareMathSymbol{\boxtimes}{\mathbin}{AMSa}{"02}

%
% Macros for moduli stacks/spaces
%
\def\QCohstack{\mathcal{QC}\!\mathit{oh}}
\def\Cohstack{\mathcal{C}\!\mathit{oh}}
\def\Spacesstack{\mathcal{S}\!\mathit{paces}}
\def\Quotfunctor{\mathrm{Quot}}
\def\Hilbfunctor{\mathrm{Hilb}}
\def\Curvesstack{\mathcal{C}\!\mathit{urves}}
\def\Polarizedstack{\mathcal{P}\!\mathit{olarized}}
\def\Complexesstack{\mathcal{C}\!\mathit{omplexes}}
% \Pic is the operator that assigns to X its picard group, usage \Pic(X)
% \Picardstack_{X/B} denotes the Picard stack of X over B
% \Picardfunctor_{X/B} denotes the Picard functor of X over B
\def\Pic{\mathop{\mathrm{Pic}}\nolimits}
\def\Picardstack{\mathcal{P}\!\mathit{ic}}
\def\Picardfunctor{\mathrm{Pic}}
\def\Deformationcategory{\mathcal{D}\!\mathit{ef}}
\usepackage{cleveref}
\crefname{originallemma}{original lemma}{original lemmas}
\Crefname{originallemma}{Original lemma}{Original lemmas}
\makeatletter
\newlabel{lemma-finite-flat-over-almost-dense-subscheme}{{0ABA}{source}{}{sourcepointer-1}{}}
\newlabel{lemma-finite-flat-over-almost-dense-subscheme@cref}{{[originallemma][0][]0ABA}{[1][0][]source}{}{}{}}
\makeatother
\crefname{originallemma}{original lemma}{original lemmas}
\Crefname{originallemma}{Original lemma}{Original lemmas}
\makeatletter
\newlabel{lemma-finite-stratify}{{0AB9}{source}{}{sourcepointer-2}{}}
\newlabel{lemma-finite-stratify@cref}{{[originallemma][0][]0AB9}{[1][0][]source}{}{}{}}
\makeatother
\crefname{originallemma}{original lemma}{original lemmas}
\Crefname{originallemma}{Original lemma}{Original lemmas}
\makeatletter
\newlabel{divisors-lemma-locally-free-rank-r-pullback}{{05P8}{source}{}{sourcepointer-3}{}}
\newlabel{divisors-lemma-locally-free-rank-r-pullback@cref}{{[originallemma][0][]05P8}{[1][0][]source}{}{}{}}
\makeatother
\crefname{originaltheorem}{original theorem}{original theorems}
\Crefname{originaltheorem}{Original theorem}{Original theorems}
\makeatletter
\newlabel{morphisms-theorem-chevalley}{{054K}{source}{}{sourcepointer-4}{}}
\newlabel{morphisms-theorem-chevalley@cref}{{[originaltheorem][0][]054K}{[1][0][]source}{}{}{}}
\makeatother
\crefname{originallemma}{original lemma}{original lemmas}
\Crefname{originallemma}{Original lemma}{Original lemmas}
\makeatletter
\newlabel{properties-lemma-generic-point-in-constructible}{{0AAW}{source}{}{sourcepointer-5}{}}
\newlabel{properties-lemma-generic-point-in-constructible@cref}{{[originallemma][0][]0AAW}{[1][0][]source}{}{}{}}
\makeatother
\title{Addendum on finite presentation: the thickening W' to W is of finite presentation, and generic-orbit points lie in W}
\author{The Stacks Project Authors}
\date{stacks-project:more-groupoids.tex:local-snapshot-20261012; GNU FDL 1.2 or later}
\begin{document}
\maketitle
\setcounter{subsection}{1}
\noindent\textit{Editorial note.} All delivered mathematical text is the author's own, copied from the pinned Stacks Project chapter file more-groupoids.tex; no statement or proof step was rewritten, reordered or removed. The counted span is the author environment and its own proof environment. References inside the span are handed to the assembler as navigation stubs that carry the canonical tag of the cited tag only; whether each of them is load-bearing is recorded on the stub itself.\par\medskip
\noindent\textit{Allowed standard prerequisites (editorial).}\begin{itemize}
\item The R-invariant thickening decomposition with the covering property (more-groupoids chapter, Lemma lemma-finite-flat-over-almost-dense-subscheme, canonical tag 0ABA; collected on this lane as L22511). cited at statement lines 2753-2754 and proof lines 2781-2787. The base lemma being strengthened: its construction W = coprod W\_r with W\_r inside Z\_\{r - 1\} without Z\_r and its covering give the frame for part (2).
\item The finite stratification U = Z\_0 superset Z\_1 superset ... (more-groupoids chapter, Lemma lemma-finite-stratify, canonical tag 0AB9). cited at proof lines 2764-2766. Provides the stratification whose finite presentation refinement and invariance drive both parts.
\item Pullback description of the strata of the finite stratification, of finite presentation when s, t are (divisors chapter, Lemma divisors-lemma-locally-free-rank-r-pullback, canonical tag 05P8). cited at proof lines 2766-2768. Gives that the closed immersions Z\_k to U are of finite presentation, from which part (1) follows.
\item Chevalley's theorem: images of constructible sets under finitely presented morphisms are constructible (morphisms chapter, Theorem morphisms-theorem-chevalley, canonical tag 054K; kind theorem). cited at proof lines 2775-2776. The images of Z\_k to U and Z\_\{k + 1\} to U are locally constructible, the input for the generic point neighbourhood argument.
\item A generic point of an irreducible component has a neighbourhood inside a constructible set containing it (properties chapter, Lemma properties-lemma-generic-point-in-constructible, canonical tag 0AAW). cited at proof lines 2777-2780. Produces the open neighbourhoods U\_i of the orbit points inside Z\_k without Z\_\{k + 1\} set theoretically.
\end{itemize}\medskip
\ifdefined\hypertarget\hypertarget{sourcepointer-1}{}\fi
\noindent\textit{Source reference.} Original lemma 0ABA: sources/stacks-chow/tags.txt, source lines 13013--13013. Cited at statement lines 2753-2754 as the base lemma being strengthened and at proof lines 2781-2787 for the construction W = coprod W\_r and the covering used in the stratum matching.\par\medskip
\ifdefined\hypertarget\hypertarget{sourcepointer-2}{}\fi
\noindent\textit{Source reference.} Original lemma 0AB9: sources/stacks-chow/tags.txt, source lines 13012--13012. Cited at proof lines 2764-2766: the stratification U = Z\_0 superset Z\_1 superset ... under finite presentation assumptions.\par\medskip
\ifdefined\hypertarget\hypertarget{sourcepointer-3}{}\fi
\noindent\textit{Source reference.} Original lemma 05P8: sources/stacks-chow/tags.txt, source lines 7332--7332. Cited at proof lines 2766-2768: the strata are closed immersions of finite presentation, yielding part (1).\par\medskip
\ifdefined\hypertarget\hypertarget{sourcepointer-4}{}\fi
\noindent\textit{Source reference.} Original theorem 054K: sources/stacks-chow/tags.txt, source lines 6643--6643. Cited at proof lines 2775-2776: the images of Z\_k to U and Z\_\{k + 1\} to U are locally constructible.\par\medskip
\ifdefined\hypertarget\hypertarget{sourcepointer-5}{}\fi
\noindent\textit{Source reference.} Original lemma 0AAW: sources/stacks-chow/tags.txt, source lines 12999--12999. Cited at proof lines 2777-2780: the open neighbourhoods U\_i of the generic orbit points inside Z\_k without Z\_\{k + 1\}.\par\medskip
\begin{lemma}
\label{lemma-finite-flat-over-almost-dense-subscheme-addendum}
In Lemma \ref{lemma-finite-flat-over-almost-dense-subscheme}
assume in addition that $s$ and $t$ are of finite presentation.
Then
\begin{enumerate}
\item the morphism $W' \to W$ is of finite presentation, and
\item if $u \in U$ is a point whose $R$-orbit consists of
generic points of irreducible components of $U$, then $u \in W$.
\end{enumerate}
\end{lemma}

\begin{proof}
In this case the stratification
$U = Z_0 \supset Z_1 \supset Z_2 \supset \ldots$ of
Lemma \ref{lemma-finite-stratify} is given by closed immersions $Z_k \to U$
of finite presentation, see
Divisors, Lemma \ref{divisors-lemma-locally-free-rank-r-pullback}.
Part (1) follows immediately from this as $W' \to W$ is locally given
by intersecting the open $W$ by $Z_r$. To see part (2)
let $\{u_1, \ldots, u_n\}$ be the orbit of $u$.
Since the closed subschemes $Z_k$ are $R$-invariant and
$\bigcap Z_k = \emptyset$, we find an $k$ such that $u_i \in Z_k$
and $u_i \not \in Z_{k + 1}$ for all $i$.
The image of $Z_k \to U$ and $Z_{k + 1} \to U$ is locally constructible
(Morphisms, Theorem \ref{morphisms-theorem-chevalley}).
Since $u_i \in U$ is a generic point of an irreducible component
of $U$, there exists an open neighbourhood $U_i$ of $u_i$ which
is contained in $Z_k \setminus Z_{k + 1}$ set theoretically
(Properties, Lemma \ref{properties-lemma-generic-point-in-constructible}).
In the proof of Lemma \ref{lemma-finite-flat-over-almost-dense-subscheme}
we have constructed $W$ as a disjoint union $\coprod W_r$
with $W_r \subset Z_{r - 1} \setminus Z_r$ such that
$U = \bigcup t(s^{-1}(\overline{W_r}))$. As $\{u_1, \ldots, u_n\}$
is an $R$-orbit we see that $u \in t(s^{-1}(\overline{W_r}))$
implies $u_i \in \overline{W_r}$ for some $i$ which implies
$U_i \cap W_r \not = \emptyset$ which implies $r = k$.
Thus we conclude that $u$ is in
$$
W_{k + 1} = Z_k \setminus
\left(
Z_{k + 1} \cup \bigcup\nolimits_{r \leq k} t(s^{-1}(\overline{W_r}))
\right)
$$
as desired.
\end{proof}

\end{document}
